% ============================================================================
%  PRECÁLCULO — Sesión 2
%  Semana 2: 17 al 22 de agosto de 2026   (lunes 17: festivo)
%  1.2 Operaciones básicas con números enteros
%  Duración prevista: 3 horas
% ============================================================================
\documentclass[aspectratio=169,10pt]{beamer}

\usepackage[utf8]{inputenc}
\IfFileExists{spanish.ldf}{\usepackage[spanish,es-noshorthands]{babel}}{}
\usepackage{amsmath,amssymb}
\usepackage{booktabs}
\usepackage{tikz}

\usetheme{Uniminuto}

\uniasignatura{Precálculo}
\uniprofesor{Andrés Felipe Puerta Velez}
\unicorreo{andres.puerta-v@uniminuto.edu.co}
\uninrc{MATE1161}
\unisemana{Sesión 2 — Semana 2 · 17 al 22 de agosto de 2026}

\newcommand{\listatiny}{\setbeamerfont{itemize/enumerate body}{size=\tiny}\tiny}

\begin{document}

\portada

\begin{frame}{Orden del día}
  \begin{enumerate}
    \item Llamado a lista y revisión de los ejercicios propuestos
    \item Repaso: dónde quedamos
    \item \azul{Tema 1.2: suma, resta, producto y cociente de enteros}
    \item Potencias de enteros y jerarquía de operaciones
    \item Divisibilidad: primos, MCD y mcm
    \item Talleres guiados, errores frecuentes y cierre
  \end{enumerate}

  \vspace{0.25cm}
  {\footnotesize\color{uniGris} Sesión de 3 horas. Lo de hoy parece elemental y
  \destacar{es donde se pierde la mitad de los exámenes} de todo el semestre:
  un signo mal puesto arruina un problema entero.}
\end{frame}

\begin{frame}{Repaso: dónde quedamos}
  \footnotesize
  En la sesión 1 vimos los conjuntos numéricos, las propiedades de las
  operaciones y el orden en $\mathbb{R}$.

  \vspace{0.2cm}
  \begin{alertblock}{La pregunta de hoy}
    \footnotesize
    Ya sabemos \azul{qué} números existen. Ahora: \destacar{¿cómo se opera con
    ellos cuando hay signos negativos de por medio?} Y sobre todo, ¿en qué
    orden se hacen las operaciones cuando aparecen todas juntas?
  \end{alertblock}

  \vspace{0.2cm}
  \begin{block}{Dos ideas de la sesión pasada que hoy se usan todo el tiempo}
    \footnotesize
    El \azul{opuesto} de $a$ es $-a$, y cumple $a + (-a) = 0$. \quad
    El \azul{valor absoluto} $|a|$ es la distancia al cero, y nunca es
    negativo.
  \end{block}
\end{frame}

\seccion{1.2 Operaciones con números enteros}

\begin{frame}{Suma y resta: dos reglas y nada más}
  \scriptsize
  \begin{columns}[T,onlytextwidth]
    \begin{column}{0.48\textwidth}
      \begin{block}{Signos iguales}
        \scriptsize
        Se \azul{suman} los valores absolutos y se conserva el signo.
        \[ -6 + (-9) = -15 \qquad 7 + 4 = 11 \]
      \end{block}
    \end{column}
    \begin{column}{0.48\textwidth}
      \begin{block}{Signos distintos}
        \scriptsize
        Se \azul{restan} los valores absolutos y se pone el signo del
        \destacar{mayor en valor absoluto}.
        \[ -8 + 15 = 7 \qquad 12 - 20 = -8 \]
      \end{block}
    \end{column}
  \end{columns}

  \vspace{0.15cm}
  \begin{exampleblock}{La resta no existe}
    \scriptsize
    Restar es sumar el opuesto: \quad $a - b = a + (-b)$
    \[ -7 - (-3) = -7 + 3 = \destacar{-4} \qquad\qquad
       5 - 9 = 5 + (-9) = \destacar{-4} \]
    Por eso \azul{menos por menos da más} en la resta: $-(-3) = +3$. Quitar una
    deuda es ganar.
  \end{exampleblock}

  \vspace{0.05cm}
  {\tiny\color{uniGris} En la recta numérica: sumar un positivo es moverse a la
  derecha, sumar un negativo es moverse a la izquierda. Nada más.}
\end{frame}

\begin{frame}{Producto y cociente: la ley de los signos}
  \scriptsize
  \begin{columns}[T,onlytextwidth]
    \begin{column}{0.44\textwidth}
      \centering
      \begin{tabular}{@{}ccc@{}}
        \toprule
        & $+$ & $-$ \\
        \midrule
        $+$ & $+$ & $-$ \\
        $-$ & $-$ & $+$ \\
        \bottomrule
      \end{tabular}

      \vspace{0.15cm}
      \[ (-5)(6) = -30 \qquad (-4)(-9) = 36 \]
      \[ \frac{-36}{-4} = 9 \qquad \frac{48}{-6} = -8 \]
    \end{column}
    \begin{column}{0.52\textwidth}
      \begin{exampleblock}{¿Por qué $(-)(-) = +$?}
        \tiny
        Porque si no, se rompería la distributiva. Observe:
        \[ (-3)\big[\,5 + (-5)\,\big] = (-3)(0) = 0 \]
        Pero distribuyendo:
        \[ (-3)(5) + (-3)(-5) = -15 + (-3)(-5) \]
        Para que el total sea $0$, necesariamente $(-3)(-5) = +15$.
      \end{exampleblock}
    \end{column}
  \end{columns}

  \vspace{0.1cm}
  \begin{alertblock}{La regla práctica para varios factores}
    \scriptsize
    Se cuentan los factores negativos: si son \destacar{pares}, el resultado es
    positivo; si son \destacar{impares}, negativo.
    \[ (-3)(4)(-2) = \destacar{+24} \qquad (-1)(-2)(-3) = \destacar{-6} \]
  \end{alertblock}
\end{frame}

\begin{frame}{Potencias de enteros: la trampa del signo}
  \scriptsize
  \begin{columns}[T,onlytextwidth]
    \begin{column}{0.48\textwidth}
      \begin{block}{Con paréntesis}
        \scriptsize
        \[ (-2)^4 = (-2)(-2)(-2)(-2) = \destacar{16} \]
        \[ (-2)^3 = (-2)(-2)(-2) = \destacar{-8} \]
        La base \azul{es} $-2$.
      \end{block}
    \end{column}
    \begin{column}{0.48\textwidth}
      \begin{alertblock}{Sin paréntesis}
        \scriptsize
        \[ -2^4 = -(2^4) = \destacar{-16} \]
        \[ -2^3 = -(2^3) = \destacar{-8} \]
        La base es $2$; el menos queda \azul{afuera}.
      \end{alertblock}
    \end{column}
  \end{columns}

  \vspace{0.15cm}
  \begin{exampleblock}{La regla que resuelve todos los casos}
    \scriptsize
    El exponente afecta \destacar{solo a lo que está pegado a él}. Si el signo
    menos no está dentro del paréntesis, no entra en la potencia.
    \[ -5^2 + (-5)^2 = -25 + 25 = \destacar{0} \]
  \end{exampleblock}

  \vspace{0.05cm}
  {\tiny\color{uniGris} Con exponente \azul{par} el paréntesis cambia todo; con
  exponente \azul{impar} los dos dan lo mismo, porque el resultado es negativo
  de cualquier forma. Verifíquelo con $-2^3$ y $(-2)^3$.}
\end{frame}

\begin{frame}{Jerarquía de operaciones}
  \scriptsize
  \begin{center}
  \begin{tabular}{@{}cp{4.4cm}p{6.0cm}@{}}
    \toprule
    & \textbf{Se hace} & \textbf{Detalle} \\
    \midrule
    1 & \azul{Paréntesis} & De adentro hacia afuera: $(\ )$, $[\ ]$, $\{\ \}$ \\
    2 & \azul{Potencias y raíces} & Antes que cualquier producto \\
    3 & \azul{Productos y cocientes} &
      \destacar{De izquierda a derecha}, en el orden en que aparecen \\
    4 & \azul{Sumas y restas} &
      \destacar{De izquierda a derecha} \\
    \bottomrule
  \end{tabular}
  \end{center}

  \vspace{0.1cm}
  \begin{alertblock}{El punto 3 es el que casi nadie respeta}
    \scriptsize
    La multiplicación \destacar{no va antes} que la división: van al mismo
    nivel y se resuelven en el orden en que aparecen.
    \[ 18 \div (-3) \cdot 2 = (-6)\cdot 2 = \destacar{-12}
       \qquad\text{y no}\qquad 18 \div (-6) = -3 \]
    Lo mismo pasa con la suma y la resta.
  \end{alertblock}
\end{frame}

\begin{frame}{Ejemplo guiado 1 — paso a paso}
  \scriptsize
  \[ 5 + 3 \cdot (8-2)^2 \div 9 - (-4) \]

  \vspace{0.05cm}
  \begin{center}
  \begin{tabular}{@{}clp{5.6cm}@{}}
    \toprule
    \textbf{Paso} & \textbf{Expresión} & \textbf{Qué se hizo} \\
    \midrule
    1 & $5 + 3\cdot(6)^2 \div 9 - (-4)$ & Paréntesis: $8-2=6$ \\
    2 & $5 + 3\cdot 36 \div 9 - (-4)$   & Potencia: $6^2 = 36$ \\
    3 & $5 + 108 \div 9 - (-4)$         & Producto (va antes por estar a la izquierda) \\
    4 & $5 + 12 - (-4)$                 & Cociente: $108\div 9 = 12$ \\
    5 & $5 + 12 + 4$                    & Resta del opuesto: $-(-4) = +4$ \\
    6 & $\destacar{21}$                 & Suma de izquierda a derecha \\
    \bottomrule
  \end{tabular}
  \end{center}

  \vspace{0.1cm}
  \begin{alertblock}{Cómo se escribe en el examen}
    \tiny
    \destacar{Una línea por paso.} No se saltan pasos ni se hacen dos
    operaciones a la vez: así es como aparecen los errores de signo, y así es
    como se pierden puntos que ya se tenían.
  \end{alertblock}
\end{frame}

\begin{frame}{Ejemplo guiado 2 — con signos de agrupación anidados}
  \scriptsize
  \[ -\big[\,2 - (5-8)\,\big] + (-3)^2 \]

  \vspace{0.1cm}
  \begin{columns}[T,onlytextwidth]
    \begin{column}{0.52\textwidth}
      \begin{block}{Desarrollo}
        \tiny
        \begin{align*}
          &= -\big[\,2 - (-3)\,\big] + (-3)^2 &&\text{paréntesis interno} \\
          &= -\big[\,2 + 3\,\big] + (-3)^2 &&\text{resta del opuesto} \\
          &= -\,[\,5\,] + 9 &&\text{corchete y potencia} \\
          &= -5 + 9 \;=\; \destacar{4}
        \end{align*}
      \end{block}
    \end{column}
    \begin{column}{0.46\textwidth}
      \begin{alertblock}{El error que se comete aquí}
        \tiny
        Escribir $-[2 - (-3)] = -2 - 3$. \destacar{El menos de afuera afecta a
        todo el corchete}, no solo al primer término: primero se resuelve lo de
        adentro y después se aplica el signo.
      \end{alertblock}
    \end{column}
  \end{columns}

  \vspace{0.1cm}
  \begin{exampleblock}{Una segunda verificación, siempre}
    \scriptsize
    Rehaga el ejercicio \azul{distribuyendo el menos} desde el principio:
    $-[2-(5-8)] = -2 + (5-8) = -2 + (-3) = -5$. Mismo resultado.
    \destacar{Dos caminos que coinciden es la mejor señal de que está bien.}
  \end{exampleblock}
\end{frame}

\begin{frame}{Ejercicio 1 — en clase}
  \begin{alertblock}{Individual (10 minutos)}
    \scriptsize
    Calcule mentalmente o en el cuaderno, en menos de un minuto cada uno:

    \vspace{0.25cm}
    \centering
    \begin{tabular}{@{}llll@{}}
      1. \ $-8 + 15$ & 2. \ $12 - 20$ & 3. \ $-6 - 9$ &
      4. \ $-7 - (-3)$ \\[0.7em]
      5. \ $(-5)(6)$ & 6. \ $(-4)(-9)$ & 7. \ $\dfrac{-36}{-4}$ &
      8. \ $\dfrac{48}{-6}$ \\[0.9em]
      9. \ $(-2)^5$ & 10. \ $-3^2$ & 11. \ $(-3)^2$ &
      12. \ $(-1)^{100}$ \\
    \end{tabular}

    \vspace{0.25cm}
    \raggedright
    Al terminar, compare con su compañero: \destacar{si difieren en alguno,
    ese es el que hay que discutir}.
  \end{alertblock}
\end{frame}

\begin{frame}{Ejercicio 1 — solución}
  \scriptsize
  \begin{center}
  \begin{tabular}{@{}llllllll@{}}
    \toprule
    1 & $7$ & 2 & $-8$ & 3 & $-15$ & 4 & $-4$ \\
    5 & $-30$ & 6 & $36$ & 7 & $9$ & 8 & $-8$ \\
    9 & $-32$ & 10 & $-9$ & 11 & $9$ & 12 & $1$ \\
    \bottomrule
  \end{tabular}
  \end{center}

  \vspace{0.15cm}
  \begin{alertblock}{Los tres que separan al grupo}
    \scriptsize
    \textbf{4.} $-7-(-3) = -7+3 = -4$. No es $-10$: el doble signo se convierte
    en suma. \quad
    \textbf{10 y 11.} $-3^2 = -9$ pero $(-3)^2 = 9$. \destacar{El paréntesis
    es la diferencia entre aprobar y no aprobar.} \quad
    \textbf{12.} $(-1)$ elevado a exponente par siempre da $1$; a exponente
    impar, $-1$.
  \end{alertblock}
\end{frame}

\begin{frame}{Ejercicio 2 — en clase}
  \begin{alertblock}{En parejas (15 minutos)}
    \scriptsize
    Resuelva mostrando \destacar{todos} los pasos, uno por línea:

    \vspace{0.25cm}
    \begin{columns}[T,onlytextwidth]
      \begin{column}{0.48\textwidth}
        \begin{enumerate}\setlength{\itemsep}{0.15cm}
          \item $20 - 3(7-9)^2$
          \item $(-2)^3 + 16 \div (-8) - 5$
          \item $-5^2 + (-5)^2$
        \end{enumerate}
      \end{column}
      \begin{column}{0.48\textwidth}
        \begin{enumerate}\setlength{\itemsep}{0.15cm}
          \setcounter{enumi}{3}
          \item $6 + \dfrac{14 - (3-8)}{-19}$
          \item $2(-3)^2 - (-4)(5)$
          \item $-\big[\,7 - (2-6)\,\big] - (-3)$
        \end{enumerate}
      \end{column}
    \end{columns}
  \end{alertblock}
\end{frame}

\begin{frame}{Ejercicio 2 — solución}
  \tiny
  \begin{columns}[T,onlytextwidth]
    \begin{column}{0.49\textwidth}
      \begin{block}{Puntos 1 a 3}
        \tiny
        \begin{align*}
          1)\;& 20 - 3(-2)^2 = 20 - 3(4) = 20-12 = \destacar{8} \\[0.25em]
          2)\;& -8 + 16\div(-8) - 5 = -8 + (-2) - 5 = \destacar{-15} \\[0.25em]
          3)\;& -25 + 25 = \destacar{0}
        \end{align*}
      \end{block}
    \end{column}
    \begin{column}{0.49\textwidth}
      \begin{block}{Puntos 4 a 6}
        \tiny
        \begin{align*}
          4)\;& 6 + \tfrac{14-(-5)}{-19} = 6 + \tfrac{19}{-19}
                = 6 - 1 = \destacar{5} \\[0.25em]
          5)\;& 2(9) - (-20) = 18 + 20 = \destacar{38} \\[0.25em]
          6)\;& -[\,7+4\,] + 3 = -11 + 3 = \destacar{-8}
        \end{align*}
      \end{block}
    \end{column}
  \end{columns}

  \vspace{0.1cm}
  \begin{alertblock}{Dos observaciones}
    \tiny
    En el \destacar{1} el $(-2)$ está elevado al cuadrado, así que da $+4$; el
    $3$ multiplica después. \quad
    En el \destacar{4} la raya de fracción \azul{agrupa}: es como un paréntesis
    invisible, y primero se resuelve todo el numerador.
  \end{alertblock}
\end{frame}

\seccion{Divisibilidad, MCD y mcm}

\begin{frame}{Múltiplos, divisores y números primos}
  \scriptsize
  \begin{columns}[T,onlytextwidth]
    \begin{column}{0.48\textwidth}
      \begin{block}{Definiciones}
        \tiny
        $b$ \azul{divide} a $a$ si existe $k$ entero con $a = b\cdot k$.
        Se escribe $b \mid a$.

        \vspace{0.08cm}
        Entonces $a$ es \azul{múltiplo} de $b$, y $b$ es \azul{divisor} de $a$.

        \vspace{0.08cm}
        Un número es \destacar{primo} si tiene exactamente dos divisores: $1$ y
        él mismo. \azul{El $1$ no es primo.}
      \end{block}
    \end{column}
    \begin{column}{0.48\textwidth}
      \begin{alertblock}{Criterios rápidos}
        \tiny
        \textbf{2:} termina en cifra par. \quad
        \textbf{3:} la suma de sus cifras es múltiplo de 3. \\
        \textbf{4:} las dos últimas cifras forman múltiplo de 4. \\
        \textbf{5:} termina en 0 o 5. \quad
        \textbf{9:} la suma de cifras es múltiplo de 9. \\
        \textbf{10:} termina en 0.
      \end{alertblock}
    \end{column}
  \end{columns}

  \vspace{0.1cm}
  \begin{exampleblock}{Descomposición en factores primos}
    \scriptsize
    Todo entero mayor que 1 se escribe de \destacar{una única forma} como
    producto de primos:
    \[ 36 = 2^2\cdot 3^2 \qquad 48 = 2^4\cdot 3 \qquad
       60 = 2^2\cdot 3\cdot 5 \]
    Este es el paso previo a todo lo demás: sin descomponer, ni el MCD ni el
    mcm salen.
  \end{exampleblock}
\end{frame}

\begin{frame}{MCD y mcm: cuál se usa cuándo}
  \scriptsize
  \begin{columns}[T,onlytextwidth]
    \begin{column}{0.48\textwidth}
      \begin{block}{Máximo común divisor}
        \tiny
        Factores \azul{comunes} con el \azul{menor} exponente.
        \[ 36 = 2^2\cdot 3^2 \qquad 48 = 2^4\cdot 3 \]
        \[ \mathrm{MCD} = 2^2\cdot 3 = \destacar{12} \]
        Sirve para \destacar{repartir en partes iguales} y para
        \destacar{simplificar fracciones}.
      \end{block}
    \end{column}
    \begin{column}{0.48\textwidth}
      \begin{block}{Mínimo común múltiplo}
        \tiny
        \azul{Todos} los factores con el \azul{mayor} exponente.
        \[ \mathrm{mcm} = 2^4\cdot 3^2 = \destacar{144} \]
        Sirve para \destacar{sincronizar ciclos} y para el
        \destacar{denominador común} de la sesión 3.
      \end{block}
    \end{column}
  \end{columns}

  \vspace{0.1cm}
  \begin{alertblock}{La verificación que nunca falla}
    \scriptsize
    Para dos números: \quad
    $\mathrm{MCD}(a,b)\cdot \mathrm{mcm}(a,b) = a\cdot b$
    \[ 12 \cdot 144 = 1\,728 = 36\cdot 48 \ \checkmark \]
    {\tiny\color{uniGris} Ojo: esta identidad vale para \azul{dos} números, no
    para tres o más.}
  \end{alertblock}
\end{frame}

\begin{frame}{Ejercicio 3 — en clase}
  \begin{alertblock}{En parejas (12 minutos)}
    \scriptsize
    \begin{enumerate}\setlength{\itemsep}{0.08cm}
      \item Descomponga en factores primos: $24$, $36$ y $60$.
      \item Calcule $\mathrm{MCD}(24,36,60)$ y $\mathrm{mcm}(24,36,60)$.
      \item Verifique que $\mathrm{MCD}(36,48)\cdot\mathrm{mcm}(36,48)
            = 36\cdot 48$.
      \item Dos rutas de bus salen de la misma terminal: una cada
            \destacar{12 minutos} y otra cada \destacar{18 minutos}. Si acaban
            de salir juntas, ¿en cuántos minutos vuelven a coincidir?
      \item Un piso rectangular de $84 \times 126$ cm se quiere cubrir con
            baldosas cuadradas \destacar{del mayor tamaño posible}, sin cortar
            ninguna. ¿Cuánto mide el lado?
    \end{enumerate}
  \end{alertblock}
\end{frame}

\begin{frame}{Ejercicio 3 — solución}
  \scriptsize
  \begin{columns}[T,onlytextwidth]
    \begin{column}{0.46\textwidth}
      \begin{block}{Puntos 1 a 3}
        \tiny
        \[ 24 = 2^3\cdot 3 \qquad 36 = 2^2\cdot 3^2 \qquad
           60 = 2^2\cdot 3\cdot 5 \]
        Comunes con menor exponente:
        \[ \mathrm{MCD} = 2^2\cdot 3 = \destacar{12} \]
        Todos con mayor exponente:
        \[ \mathrm{mcm} = 2^3\cdot 3^2\cdot 5 = \destacar{360} \]
        \[ 12\cdot 144 = 1\,728 = 36\cdot 48 \ \checkmark \]
      \end{block}
    \end{column}
    \begin{column}{0.50\textwidth}
      \begin{exampleblock}{Puntos 4 y 5}
        \tiny
        \textbf{4.} Coinciden en un múltiplo común de 12 y 18, y el primero es
        el \azul{mcm}:
        \[ 12 = 2^2\cdot 3,\quad 18 = 2\cdot 3^2 \Rightarrow
           \mathrm{mcm} = 2^2\cdot 3^2 = \destacar{36 \text{ min}} \]

        \textbf{5.} El lado debe dividir a 84 y a 126, y ser el mayor posible:
        es el \azul{MCD}.
        \[ 84 = 2^2\cdot 3\cdot 7,\quad 126 = 2\cdot 3^2\cdot 7 \Rightarrow
           \mathrm{MCD} = 2\cdot 3\cdot 7 = \destacar{42 \text{ cm}} \]
      \end{exampleblock}
    \end{column}
  \end{columns}

  \vspace{0.05cm}
  {\tiny\color{uniGris} \destacar{Cómo saber cuál usar:} si el problema
  \azul{parte o reparte} algo, es MCD; si algo \azul{se repite y se
  reencuentra}, es mcm.}
\end{frame}

\begin{frame}{Ejercicio 4 — en clase (taller principal)}
  \begin{alertblock}{En grupos de 3 (20 minutos)}
    \scriptsize
    \begin{enumerate}\setlength{\itemsep}{0.06cm}
      \item En Bogotá la temperatura era de $-8\,^\circ$C a las 5 a.m., subió
            $15\,^\circ$ al mediodía y bajó $6\,^\circ$ en la noche.
            ¿Cuál fue la temperatura final?
      \item Una cuenta tiene un sobregiro de \$120\,000. Se consigna
            \$450\,000 y luego se pagan \$85\,000. ¿Cuál es el saldo?
      \item Un servidor registra estas variaciones de memoria libre, en MB:
            $-340$, $+1\,200$, $-560$, $-90$, $+75$. ¿Cuál es la variación
            neta?
      \item Escriba una sola expresión con paréntesis, potencias, producto y
            división cuyo resultado sea exactamente $\destacar{10}$; después
            intercámbiela con otro grupo para que la resuelva.
      \item Encuentre el error: \\
            \hspace*{0.3cm} $-4^2 + (-6)\div 2 \cdot 3 = 16 - 6 \div 6 = 16 - 1 = 15$
    \end{enumerate}
  \end{alertblock}
\end{frame}

\begin{frame}{Ejercicio 4 — solución}
  \scriptsize
  \begin{columns}[T,onlytextwidth]
    \begin{column}{0.46\textwidth}
      \begin{block}{Puntos 1 a 3}
        \tiny
        \textbf{1.} $-8 + 15 - 6 = \destacar{1\,^\circ\text{C}}$

        \vspace{0.1cm}
        \textbf{2.} $-120\,000 + 450\,000 - 85\,000
        = \destacar{\$245\,000}$

        \vspace{0.1cm}
        \textbf{3.} $-340 + 1\,200 - 560 - 90 + 75 = \destacar{+285}$ MB:
        la memoria libre \azul{aumentó}.
      \end{block}
    \end{column}
    \begin{column}{0.50\textwidth}
      \begin{alertblock}{Punto 5: los dos errores}
        \tiny
        \textbf{(a)} $-4^2 = \destacar{-16}$, no $+16$: el menos está fuera de
        la potencia.

        \vspace{0.08cm}
        \textbf{(b)} $(-6)\div 2\cdot 3$ va \destacar{de izquierda a derecha}:
        $(-3)\cdot 3 = -9$, no $-6\div 6$.

        \vspace{0.08cm}
        Lo correcto: \ $-16 + (-9) = \destacar{-25}$.
      \end{alertblock}
    \end{column}
  \end{columns}

  \vspace{0.05cm}
  \begin{exampleblock}{Punto 4: un ejemplo posible}
    \tiny
    \[ (3-1)^3 + 12 \div (-6) \cdot (-1) = 8 + 2 = 10 \]
    Hay infinitas respuestas correctas. Lo que se evalúa es que
    \destacar{use las cuatro cosas} y que al resolverla dé 10.
  \end{exampleblock}
\end{frame}

\begin{frame}{Errores frecuentes en esta sesión}
  \scriptsize
  \begin{center}
  \begin{tabular}{@{}p{5.2cm}p{6.6cm}@{}}
    \toprule
    \textbf{Error} & \textbf{Corrección} \\
    \midrule
    $-3^2 = 9$ &
      $-3^2 = -9$; solo $(-3)^2 = 9$ \\
    \addlinespace
    $-7-(-3) = -10$ &
      $-7+3 = -4$: el doble signo se vuelve suma \\
    \addlinespace
    Hacer la multiplicación antes que la división &
      Van al mismo nivel: de \azul{izquierda a derecha} \\
    \addlinespace
    $-[2-(-3)] = -2-3$ &
      El menos afecta a \azul{todo} el corchete: $-[5] = -5$ \\
    \addlinespace
    Operar el numerador y el denominador por separado sin agrupar &
      La raya de fracción es un paréntesis invisible \\
    \addlinespace
    Confundir MCD con mcm &
      MCD: comunes con menor exponente. mcm: todos con mayor exponente \\
    \bottomrule
  \end{tabular}
  \end{center}
\end{frame}

\begin{frame}{Cierre de la sesión}
  \begin{block}{Lo que debe quedar claro hoy}
    \scriptsize
    \begin{itemize}\setlength{\itemsep}{0.02cm}
      \item Signos iguales se suman; signos distintos se restan y manda el
            mayor en valor absoluto.
      \item Restar es \azul{sumar el opuesto}: $a - b = a + (-b)$.
      \item Ley de signos: pares de negativos dan positivo, impares dan
            negativo.
      \item $-a^n \neq (-a)^n$ cuando $n$ es \destacar{par}.
      \item Jerarquía: paréntesis, potencias, $\times\div$ y $+-$, estos dos
            últimos \azul{de izquierda a derecha}.
      \item MCD para repartir; mcm para sincronizar y para denominadores
            comunes.
    \end{itemize}
  \end{block}

  \vspace{0.08cm}
  \begin{alertblock}{Trabajo independiente}
    \tiny
    Resolver los 20 ejercicios de la lámina siguiente. En la sesión 3 pasamos a
    los \destacar{racionales}: sumar, restar, multiplicar y dividir fracciones,
    donde el mcm de hoy se vuelve la herramienta principal.
  \end{alertblock}
\end{frame}

\begin{frame}{Ejercicios propuestos}
  \vspace{-0.1cm}
  \begin{beamercolorbox}[wd=\textwidth,sep=0.10cm]{block body}
    \centering\scriptsize
    \textbf{1--8:} signos. \quad \textbf{9--15:} jerarquía. \quad
    \textbf{16--20:} divisibilidad y aplicaciones.
  \end{beamercolorbox}
  \listatiny
  \vspace{0.15cm}
  \begin{columns}[T,onlytextwidth]
    \begin{column}{0.47\textwidth}
      \begin{enumerate}\setlength{\itemsep}{0.16cm}
        \item $-14 + 23$
        \item $-11 - 17$
        \item $9 - (-15)$
        \item $(-7)(-8)$
        \item $(-72) \div 9$
        \item $(-2)(3)(-4)(-1)$
        \item $(-1)^{37}$
        \item $-4^3$ \ y \ $(-4)^3$
        \item $15 - 4 \cdot 3$
        \item $(15-4)\cdot 3$
      \end{enumerate}
    \end{column}
    \begin{column}{0.50\textwidth}
      \begin{enumerate}\setlength{\itemsep}{0.16cm}
        \setcounter{enumi}{10}
        \item $24 \div 6 \cdot 2$
        \item $-3^2 + (-3)^2 - 3^2$
        \item $8 - [\,5 - (2-7)\,]$
        \item $\dfrac{(-4)^2 - 6}{-5}$
        \item $2 - 3[\,4 - (6-9)\,]$
        \item Descomponga $180$ en factores primos
        \item $\mathrm{MCD}(45, 75)$
        \item $\mathrm{mcm}(8, 12, 20)$
        \item Tres luces parpadean cada 6, 10 y 15 s. ¿Cada cuánto coinciden?
        \item ¿Es $2\,346$ divisible por 3? ¿Y por 4?
      \end{enumerate}
    \end{column}
  \end{columns}
\end{frame}

\begin{frame}{Respuestas}{Autoevaluación de los ejercicios propuestos}
  \listatiny
  \vspace{0.15cm}
  \begin{columns}[T,onlytextwidth]
    \begin{column}{0.47\textwidth}
      \begin{enumerate}\setlength{\itemsep}{0.16cm}
        \item $9$
        \item $-28$
        \item $24$
        \item $56$
        \item $-8$
        \item $-24$ \ (tres negativos: impar)
        \item $-1$
        \item $-64$ \ y \ $-64$ \ (exponente impar: coinciden)
        \item $15-12 = 3$
        \item $11\cdot 3 = 33$
      \end{enumerate}
    \end{column}
    \begin{column}{0.50\textwidth}
      \begin{enumerate}\setlength{\itemsep}{0.16cm}
        \setcounter{enumi}{10}
        \item $4\cdot 2 = 8$
        \item $-9+9-9 = -9$
        \item $8 - [\,5+5\,] = -2$
        \item $\dfrac{16-6}{-5} = -2$
        \item $2 - 3[\,7\,] = -19$
        \item $180 = 2^2\cdot 3^2\cdot 5$
        \item $15$
        \item $120$
        \item $\mathrm{mcm}(6,10,15) = 30$ s
        \item Por 3 sí ($2+3+4+6=15$); por 4 no ($46$ no es múltiplo de 4)
      \end{enumerate}
    \end{column}
  \end{columns}
\end{frame}

\begin{frame}{Referencias}
  \footnotesize
  \begin{itemize}
    \item[{[1]}] Ramos del Olmo, Á. M. (2016). \textit{Matemáticas básicas
      para el acceso a la universidad}. Ediciones Pirámide.
      \texttt{https://elibro.net/es/ereader/uniminuto/49134}
    \item[{[2]}] González Pulido, J. W., Chacón Benavides, J. A. \&
      Fonseca Correa, L. Á. (2023). \textit{Matemática básica} (1.\textsuperscript{a} ed.).
      Universidad Mariana.
      \texttt{https://elibro.net/es/ereader/uniminuto/280647}
    \item[{[3]}] Placencia Valero, J. (2008). \textit{Compendio de matemática
      básica elemental} (1.\textsuperscript{a} ed.). Editorial Tébar Flores.
      \texttt{https://elibro.net/es/ereader/uniminuto/269758}
  \end{itemize}

  \vspace{0.35cm}
  \begin{center}
    {\Large\color{uniAzul}\textbf{¿Preguntas?}}
  \end{center}
\end{frame}

\end{document}
